\documentclass[11pt,a4paper]{article}
% Common preamble for the hint sheets (\input from each file)
\usepackage[utf8]{inputenc}
\usepackage[T1]{fontenc}
\usepackage{lmodern}
\usepackage{amsmath,amssymb,amsthm,mathtools}
\usepackage[a4paper,margin=2.3cm]{geometry}
\usepackage{enumitem}
\usepackage{xcolor}
\usepackage[hidelinks]{hyperref}
\renewcommand{\leq}{\leqslant}
\renewcommand{\le}{\leqslant}
\renewcommand{\geq}{\geqslant}
\renewcommand{\ge}{\geqslant}
\setlength{\parindent}{0pt}
\setlength{\parskip}{0.35em}
% \hintstitle{Sheet title}{Class work / Homework}
\newcommand{\hintstitle}[2]{\begin{center}{\LARGE\bfseries Hints: #1\par}\vspace{0.5em}{\large #2\par}\end{center}
\noindent\textit{The hints to each problem are cumulative. Read one, try again, and open the next only if you are still stuck.}\par\medskip\hrule\medskip}
\newcommand{\problem}[1]{\par\bigskip{\large\bfseries Problem #1}\par\nopagebreak\smallskip}
\newcommand{\subpart}[1]{\par\smallskip\textbf{(#1)}\par\nopagebreak}
\newcommand{\hint}[1]{\par\smallskip\noindent\textcolor{black!70}{\textbf{Hint #1.}}\enspace}
\newcommand{\note}[1]{\par\smallskip\noindent\textit{Note.} #1\par}
\providecommand{\lt}{<}
\providecommand{\gt}{>}
\providecommand{\tightlist}{\setlength{\itemsep}{0pt}\setlength{\parskip}{0pt}}
\emergencystretch=3em

\begin{document}
\hintstitle{Limits -- 6. Asymptotics and sharp estimates: sums, integrals, Stirling}{Class work}

\problem{1}
\hint{1} For monotone $f$ compare $f(k)$ with $\int_k^{k+1}f$ on each unit interval: for increasing $f$ one has $f(k)\le\int_k^{k+1}f\le f(k+1)$. Summing gives $\sum_{k=1}^{n-1}f(k)\le\int_1^n f\le\sum_{k=2}^n f(k)$, so the difference $\sum_{k=1}^n f(k)-\int_1^n f$ lies between $f(1)$ and $f(n)$.
\hint{2} For decreasing $f\ge 0$ put $d_n=\sum_{k\le n}f(k)-\int_1^n f$. Then $d_n-d_{n+1}=\int_n^{n+1}f-f(n+1)\ge 0$ and $d_n\ge f(n)\ge 0$, so $d_n$ converges. With $f(x)=1/x$ this gives $\gamma=\lim(H_n-\ln n)$.
\hint{3} For the next term write $H_n-\ln n-\gamma=\sum_{k\ge n}\bigl(\ln(1+\tfrac1k)-\tfrac1{k+1}\bigr)$ and expand: the summand is $\frac1{2k^2}+O(k^{-3})$. Since $\sum_{k\ge n}k^{-2}=\frac1n+O(n^{-2})$, you get $H_n=\ln n+\gamma+\frac1{2n}+O(n^{-2})$.
\hint{4} Numerically: $H_{10}=2.928968\ldots$, $\ln 10+\gamma=2.879801\ldots$. The two-term approximation $\ln n+\gamma$ is off by about $0.049\approx\frac1{2n}$; adding $\frac1{20}$ leaves an error of about $-0.00083\approx-\frac1{12n^2}$.

\problem{2}
\hint{1} Take $f=\ln x$ and compare $\ln n!=\sum\ln k$ with the trapezoid rule rather than with the bare integral: consider $d_n=\ln n!-(n+\tfrac12)\ln n+n$.
\hint{2} Compute $d_n-d_{n+1}=(n+\tfrac12)\ln\frac{n+1}{n}-1$. With $t=\frac1{2n+1}$ one has $\frac{n+1}{n}=\frac{1+t}{1-t}$, so $d_n-d_{n+1}=\frac{1}{2t}\ln\frac{1+t}{1-t}-1=\frac{t^2}3+\frac{t^4}5+\dots$
\hint{3} From this, $0<d_n-d_{n+1}<\frac13\sum_{j\ge1}t^{2j}=\frac1{12}\bigl(\frac1n-\frac1{n+1}\bigr)$. So $d_n$ decreases and $d_n-\frac1{12n}$ increases; both converge to some $C$. This gives $n!=e^{C}\sqrt n\,(n/e)^n e^{\theta_n}$ with $0<\theta_n<\frac1{12n}$. (The lower bound $\frac1{12n+1}$ needs a slightly sharper estimate of the series from below.)
\hint{4} For the constant use Wallis: with $a_m=\int_0^{\pi/2}\sin^m x\,dx$ one has $a_{2n}/a_{2n+1}\to1$, which gives $\frac{\pi}{2}=\lim\frac{2^{4n}(n!)^4}{((2n)!)^2(2n+1)}$. Substitute $n!\sim e^C\sqrt n(n/e)^n$ and solve to get $e^{C}=\sqrt{2\pi}$.

\problem{3}
\hint{1} Write $S_n=\frac1n\sum_{k=1}^n f(k/n)$ with $f(x)=\frac1{\sqrt{1+x^2}}$. This is a Riemann sum, and $\int_0^1 f=\operatorname{arsinh}1=\ln(1+\sqrt2)$.
\hint{2} For the second limit you need the error of the right-endpoint Riemann sum of a smooth $f$. On $[\frac{k-1}n,\frac kn]$, expanding $f$ around the right endpoint gives $\frac1nf(\frac kn)-\int f=\frac{1}{2n^2}f'(\frac kn)+O(n^{-3})$. Summing gives $\frac{f(1)-f(0)}{2n}+O(n^{-2})$.
\hint{3} So $n\bigl(S_n-\ln(1+\sqrt2)\bigr)\to\frac{f(1)-f(0)}2=\frac{1}{2}\bigl(\frac1{\sqrt2}-1\bigr)=\frac{\sqrt2-2}{4}$.

\problem{4}
\hint{1} Let $I_m=\int_0^1\frac{dx}{1+x^m}$. Show that $\frac12<I_m<1$. Then $c_m=(1/I_m-1)^{1/m}$ is forced, and it lies in $(0,1)$.
\hint{2} You need $kc_k^k=k\,\frac{1-I_k}{I_k}$. Write $1-I_k=\int_0^1\frac{x^k}{1+x^k}dx$ and substitute $t=x^k$.
\hint{3} You get $1-I_k=\frac1k\int_0^1\frac{t^{1/k}}{1+t}\,dt$, and the integral tends to $\int_0^1\frac{dt}{1+t}=\ln2$, while $I_k\to1$. The limit is $\ln 2$.

\problem{5}
\hint{1} Use the first step of Euler--Maclaurin from Theorem 2: $\sum_{k=1}^n f(k)=\int_1^n f+\frac{f(1)+f(n)}2+R_n$ with $R_n=\int_1^n\frac{\{x\}(1-\{x\})}2f''(x)\,dx$. If $\int_1^\infty|f''|<\infty$, then $R_n=R_\infty+o(1)$, and everything constant goes into one constant $C$.
\hint{2} For $f=\sqrt x$: $f''\sim x^{-3/2}$ is integrable, so $\sum_{k\le n}\sqrt k=\frac23n^{3/2}+\frac12\sqrt n+C+o(1)$. The terms used are the integral, the boundary term $f(n)/2$ and the convergent remainder. The next term, $\frac{B_2}{2}f'(n)=\frac1{24\sqrt n}$, is already $o(1)$. (In fact $C=\zeta(-\frac12)\approx-0.2079$.)
\hint{3} For $f=\ln^2x$: $\int_1^n\ln^2x\,dx=n\ln^2n-2n\ln n+2n-2$, and $f''=\frac{2-2\ln x}{x^2}$ is integrable. So $\sum_{k\le n}\ln^2k=n\ln^2n-2n\ln n+2n+\frac12\ln^2n+C'+o(1)$. (In fact $C'=\zeta''(0)\approx-2.0064$.)

\end{document}
